\documentclass[11pt]{article}
\usepackage[margin=1in]{geometry}
\usepackage{amsmath,amssymb,amsthm}
\usepackage{graphicx}
\usepackage{microtype}
\usepackage{xcolor}
\usepackage{enumitem}
\usepackage[colorlinks=true,linkcolor=blue!55!black,citecolor=blue!55!black,urlcolor=blue!55!black]{hyperref}

\newtheorem{theorem}{Theorem}
\newtheorem{lemma}[theorem]{Lemma}
\newtheorem{corollary}[theorem]{Corollary}
\newcommand{\R}{\mathbb R}
\newcommand{\N}{\mathbb N}
\newcommand{\dist}{\operatorname{dist}}
\newcommand{\rinf}{\mathrm{rinf}}

\title{Intrinsic and extrinsic reinforced Triebel--Lizorkin spaces on a cube\
\large Candidate full solution of Conjecture 4.22 in arXiv:1302.3751}
\author{Automated Mathematics Run \texttt{fa\_banach\_001}, lane 12}
\date{9 August 2026}

\begin{document}
\maketitle

\begin{center}
\fcolorbox{blue!55!black}{blue!4}{\parbox{0.89\textwidth}{
\textbf{Status.} Candidate full solution, subject to human verification.
For the entire parameter range in Conjecture 4.22,
\[
 F_{p,q}^{s,\rinf}(Q)=F_{p,q}^{s,\rinf}(Q)^*
\]
with equivalent quasi-norms.  More strongly, the proof constructs a bounded
linear right inverse to restriction which preserves all reinforced face
conditions simultaneously.}}
\end{center}

\section{Source and exact conjecture}

The source is B.~Scharf, \emph{Wavelets in function spaces on cellular
domains}~\cite{Scharf}, arXiv:1302.3751.  Let $Q=(0,1)^n$, $n\ge2$.
For every open $l$-dimensional face $\Gamma$ of $Q$, put
$k=n-l$ and $d_\Gamma(x)=\dist(x,\Gamma)$.  A face is critical when
\[
 r_\Gamma=s-\frac{k}{p}\in\N_0.
\]
For such a face the reinforced condition requires, for every multi-index
$\alpha$ normal to $\Gamma$ with $|\alpha|=r_\Gamma$,
\begin{equation}\label{eq:reinforced}
 d_\Gamma^{-k/p}D^\alpha f\in L_p(U_\Gamma),
 \qquad U_\Gamma=\{x:d_\Gamma(x)<\varepsilon\}.
\end{equation}

The extrinsic space $F_{p,q}^{s,\rinf}(Q)$ consists of restrictions to $Q$
of global $F_{p,q}^s(\R^n)$ functions satisfying \eqref{eq:reinforced} on
both sides of every critical face.  The intrinsic space
$F_{p,q}^{s,\rinf}(Q)^*$ asks for the same weighted estimates only on
$U_\Gamma\cap Q$.  Restriction immediately gives
\[
 F_{p,q}^{s,\rinf}(Q)\hookrightarrow F_{p,q}^{s,\rinf}(Q)^*.
\]

Conjecture 4.22, on source PDF page 96, asserts that this embedding is an
equality whenever
\[
 1\le p<\infty,\qquad 0<q<\infty,\qquad s>\sigma_{p,q}.
\]
The source explicitly proposes constructing an extension operator which
preserves the reinforced conditions.  Figure~\ref{fig:conjecture} records
the conjecture and this proposed route.

\begin{figure}[ht]
 \centering
 \includegraphics[width=\textwidth]{figures/open_problem_crop.png}
 \caption{Conjecture 4.22 and its extension-operator route on page 96 of
 the arXiv source~\cite{Scharf}.}
 \label{fig:conjecture}
\end{figure}

We use the intrinsic quasi-norm
\begin{equation}\label{eq:star-norm}
 \|f\|_*=\|f\|_{F_{p,q}^s(Q)}+
 \sum_{\Gamma\ \mathrm{critical}}\ \sum_{\substack{\alpha\perp\Gamma\\
 |\alpha|=r_\Gamma}}
 \left(\int_{U_\Gamma\cap Q}|D^\alpha f(x)|^p
 \frac{dx}{d_\Gamma(x)^k}\right)^{1/p}.
\end{equation}
Changing the sufficiently small common value of $\varepsilon$ changes this
quasi-norm only by an equivalent one.  Indeed, away from a smaller face
collar the weight is bounded, while
$D^\alpha f\in F_{p,q}^{s-|\alpha|}(Q)\subset L_p(Q)$ because
$s-|\alpha|=k/p>0$.

\section{Proof intuition}

The cube admits an extension which, in a small collar, is a finite linear
combination of coordinate reflections and dilations back into $Q$.  A
finite common Seeley--Hestenes reflection can be chosen bounded on the one
fixed space $F_{p,q}^s$; no infinite reflection series is needed.
Tensoring the construction over the coordinates gives, on every collar
cell,
\[
 D^\alpha E_Qf(x)=\sum_A c_{A,\alpha}(D^\alpha f)(T_Ax),
\]
where each $T_A$ is diagonal affine and maps the cell into $Q$.

The decisive geometric fact is one-sided in exactly the useful direction.
If $x$ is close to an open face $\Gamma$, then
\[
 d_\Gamma(T_Ax)\le C,d_\Gamma(x).
\]
Normal distances are merely multiplied by one of finitely many slopes.
Tangential coordinates which lie outside the open face interval are
reflected into $(0,1)$, so their contribution to the distance disappears.
Consequently
$d_\Gamma(x)^{-k}\le C^k d_\Gamma(T_Ax)^{-k}$.  A change of variables
therefore bounds the full two-sided weighted norm of the extension by the
one-sided intrinsic norm.  The same finite operator works for all faces,
dimensions, and critical derivative orders at once.

\section{A finite reflection extension on the cube}

We isolate the standard extension input in the form actually used below.

\begin{lemma}[Fixed-parameter finite reflection extension]\label{lem:extension}
Fix $n\ge1$, $0<p<\infty$, $0<q<\infty$, and $s\in\R$.  There is a bounded
linear right inverse
\[
 E_Q:F_{p,q}^s(Q)\longrightarrow F_{p,q}^s(\R^n)
\]
to restriction with the following additional property.  There are
$\rho>0$, a finite set of nonzero coefficients, and a finite set of
positive slopes such that, on each cell of the $\rho$-collar of $\bar Q$,
$E_Qf$ is a finite sum
\begin{equation}\label{eq:affine-sum}
 E_Qf(x)=\sum_A c_A f(T_Ax).
\end{equation}
Every $T_A$ is a diagonal affine map from that cell into $Q$; in each
coordinate it is the identity or a reflection/dilation about $0$ or $1$.
The operator may be chosen compactly supported and its cutoff is identically
one on the whole $\rho$-collar.
\end{lemma}

\begin{proof}
We recall why the usual common reflection extension has precisely this
form.  The quantitative common-extension theorem of Lu and Yao
\cite[Theorem~23]{LuYao} states that, after choosing sufficiently large
integers $m_1,m_2$, a half-space reflection
\begin{equation}\label{eq:half-reflection}
 Rh(x',t)=\sum_{\nu=1}^N a_\nu h(x',-b_\nu t),
 \qquad t<0,\qquad b_\nu>0,
\end{equation}
is bounded from $F_{p,q}^s(\R^n_+)$ to $F_{p,q}^s(\R^n)$ provided its
finite coefficients satisfy
\[
 \sum_{\nu=1}^N a_\nu(-b_\nu)^j=1
 \quad(-m_1\le j\le m_2).
\]
Finite families satisfying as many of these moments as desired are given
explicitly in \cite[(29)--(30)]{LuYao}.  We choose the moment range larger
than every derivative order used below.

Translate and reflect \eqref{eq:half-reflection} to treat both endpoints
of $(0,1)$.  A smooth partition of unity consisting of a left endpoint
piece, an interior piece, and a right endpoint piece localizes the
half-space theorem; the interior piece is extended by zero.  The endpoint
cutoffs are constant on smaller collars.  This produces a bounded interval
extension which, on those collars, is exactly the appropriate finite
reflection formula.  To spell out the product step, set
$D_j=\R^j\times(0,1)^{n-j}$ for $0\le j\le n$.  Extend from $D_{j-1}$ to
$D_j$ by applying the interval formula in coordinate $j$.  If $g$ is any
global representative of an element of $F_{p,q}^s(D_{j-1})$, the localized
half-space theorem gives a global representative of this new function with
norm at most $C\|g\|_{F_{p,q}^s(\R^n)}$.  Taking the infimum over $g$ proves
boundedness $F_{p,q}^s(D_{j-1})\to F_{p,q}^s(D_j)$.  Composing these $n$
maps gives the asserted cube extension.

Expanding the finite tensor product yields \eqref{eq:affine-sum}.  Finally,
multiply by a smooth compactly supported function which is one on a still
smaller collar of $\bar Q$.  Smooth multiplication is bounded on
$F_{p,q}^s$, and all claimed properties remain valid after decreasing
$\rho$.
\end{proof}

For a face it is useful to make its distance completely explicit.  Write
\begin{equation}\label{eq:face}
 \Gamma=\Gamma_{I,\epsilon}
 =\{x\in\R^n:x_i=\epsilon_i\ (i\in I),\ 0<x_j<1\ (j\notin I)\},
\end{equation}
where $\epsilon_i\in\{0,1\}$ and $|I|=k$.  If
$\delta(t)=\dist(t,(0,1))$, then
\begin{equation}\label{eq:distance-formula}
 d_\Gamma(x)^2=
 \sum_{i\in I}|x_i-\epsilon_i|^2+
 \sum_{j\notin I}\delta(x_j)^2.
\end{equation}

\begin{lemma}[Every face weight is reflection-stable]\label{lem:distance}
After decreasing $\rho$ in Lemma~\ref{lem:extension}, there is $B\ge1$
such that for every open face $\Gamma$, every collar cell meeting
$\{d_\Gamma<\rho\}$, every affine term $T_A$ on that cell, and every point
$x$ in the intersection,
\begin{equation}\label{eq:distance-comparison}
 d_\Gamma(T_Ax)\le B d_\Gamma(x).
\end{equation}
The constant is independent of the face and of the cell.
\end{lemma}

\begin{proof}
Take $\rho<1/4$ and smaller, if necessary, so that every reflected point
in the construction lies in $Q$.  For $i\in I$, closeness to $\Gamma$
forces $x_i$ to be close to the specified endpoint $\epsilon_i$, not the
opposite endpoint.  The corresponding component of $T_A$ is either the
identity or a reflection/dilation about $\epsilon_i$, and hence
\[
 |(T_Ax)_i-\epsilon_i|\le B|x_i-\epsilon_i|
\]
for the maximum $B$ of the finitely many slopes.  For $j\notin I$,
$(T_Ax)_j\in(0,1)$, so $\delta((T_Ax)_j)=0$.  Formula
\eqref{eq:distance-formula} proves \eqref{eq:distance-comparison}.
\end{proof}

\section{The conjecture}

\begin{theorem}[Conjecture 4.22]\label{thm:main}
Let $Q=(0,1)^n$, $n\ge2$, and suppose
\[
 1\le p<\infty,\qquad 0<q<\infty,\qquad s>\sigma_{p,q}.
\]
Then
\[
 F_{p,q}^{s,\rinf}(Q)=F_{p,q}^{s,\rinf}(Q)^*
\]
with equivalent quasi-norms.  In fact the operator $E_Q$ in
Lemma~\ref{lem:extension} restricts to a bounded linear map
\[
 E_Q:F_{p,q}^{s,\rinf}(Q)^*\longrightarrow
 F_{p,q}^{s,\rinf}(\R^n\setminus\partial Q)
\]
and remains a right inverse to restriction.
\end{theorem}

\begin{proof}
Let $f\in F_{p,q}^{s,\rinf}(Q)^*$ and set $F=E_Qf$.  Its ordinary
$F_{p,q}^s(\R^n)$ norm is at most a constant times
$\|f\|_{F_{p,q}^s(Q)}$ by Lemma~\ref{lem:extension}.  Fix a critical face
$\Gamma=\Gamma_{I,\epsilon}$, write $k=|I|$, and let
$r=s-k/p\in\N_0$.  Fix a normal multi-index $\alpha$, so
$\alpha_j=0$ for $j\notin I$ and $|\alpha|=r$.

Choose the reinforced collar inside the region where the cutoff from
Lemma~\ref{lem:extension} is one.  On the interior of each collar cell,
differentiating the finite affine formula gives
\begin{equation}\label{eq:derivative-sum}
 D^\alpha F(x)=\sum_A c_{A,\alpha}(D^\alpha f)(T_Ax).
\end{equation}
This also holds distributionally across cell boundaries; alternatively,
those boundaries have measure zero for the weighted $L_p$ estimate.  The
moment order in Lemma~\ref{lem:extension} was chosen larger than $r$.

For a collar cell $C$, the finite-sum inequality (using $p\ge1$), followed
by a change of variables $y=T_Ax$, gives
\begin{align*}
 &\int_{C\cap U_\Gamma}|D^\alpha F(x)|^p
          \frac{dx}{d_\Gamma(x)^k}\\
 &\quad\le C\sum_A\int_{C\cap U_\Gamma}
       |D^\alpha f(T_Ax)|^p\frac{dx}{d_\Gamma(x)^k}\\
 &\quad\le C\sum_A\int_{T_A(C\cap U_\Gamma)}
       |D^\alpha f(y)|^p\frac{dy}{d_\Gamma(y)^k}.
\end{align*}
For the second inequality we used
\eqref{eq:distance-comparison} in the form
$d_\Gamma(x)^{-k}\le B^k d_\Gamma(T_Ax)^{-k}$, together with the uniformly
bounded Jacobians of the finitely many affine maps.  Moreover
$T_A(C\cap U_\Gamma)$ lies in $Q\cap\{d_\Gamma<B\varepsilon\}$.

There are only finitely many cells and affine terms.  Summing, and using
the independence of the intrinsic condition from the precise small collar
width, yields
\begin{equation}\label{eq:weighted-bound}
 \left(\int_{U_\Gamma}|D^\alpha F(x)|^p
       \frac{dx}{d_\Gamma(x)^k}\right)^{1/p}
 \le C\left(\int_{Q\cap U'_\Gamma}|D^\alpha f(y)|^p
       \frac{dy}{d_\Gamma(y)^k}\right)^{1/p}.
\end{equation}
The construction and estimate apply simultaneously to every critical face
and every required normal multi-index.  Thus
\[
 \|E_Qf\|_{F_{p,q}^{s,\rinf}(\R^n\setminus\partial Q)}
 \le C\|f\|_*.
\]
Because $E_Qf|_Q=f$, this proves
$F_{p,q}^{s,\rinf}(Q)^*\hookrightarrow F_{p,q}^{s,\rinf}(Q)$.
The reverse embedding is the immediate restriction observation from the
source, and the displayed estimates give equivalence of quasi-norms.
\end{proof}

\begin{corollary}[The intrinsic Sobolev example]\label{cor:sobolev}
For the unit square $Q\subset\R^2$,
\[
 W_2^{1,\rinf}(Q)^*
 =\left\{f\in W^{1,2}(Q):
   \int_Q\frac{|f(x)|^2}{\dist(x,\Gamma_0)^2}\,dx<\infty\right\}
 =F_{2,2}^{1,\rinf}(Q),
\]
where $\Gamma_0$ is the set of four vertices, with equivalent norms.
Consequently the intrinsic space has an oscillating $u$-Riesz basis for
every integer $u>1$.
\end{corollary}

\begin{proof}
For $n=2$, $s=1$, and $p=q=2$, the only critical stratum has $l=0$,
$k=2$, and $r=0$, so \eqref{eq:reinforced} is exactly the displayed corner
condition.  Theorem~\ref{thm:main} identifies the intrinsic and extrinsic
spaces.  The source's wavelet theorem for the latter therefore transfers
to the former.
\end{proof}

\section{Verification, scope, and novelty status}

The proof is analytic; no numerical experiment is needed.  The following
checks are explicit in the argument.
\begin{itemize}[leftmargin=*]
 \item The external functional-space input is precisely the finite
 common-extension case of Lu--Yao Theorem~23, not their infinite universal
 reflection series.  Finiteness is what makes the weighted argument
 immediate.
 \item Formula \eqref{eq:distance-formula} includes the tangential endpoint
 contribution to the distance from an \emph{open} face.  Under reflection
 that contribution becomes zero; it is not silently discarded.
 \item The cutoff is one on all reinforced collars, so no derivative of the
 cutoff appears in \eqref{eq:derivative-sum}.
 \item The restriction $p\ge1$ is used exactly in the finite-sum $L_p$
 estimate.  The proof covers every $0<q<\infty$ in the conjecture.
\end{itemize}

The result is for an axis-parallel cube.  Extending the same proof to a
general polyhedron would require a separate construction compatible with
intersecting non-orthogonal faces and is not claimed here.

Targeted searches through 9 August 2026 for the arXiv id, the exact
intrinsic/extrinsic reinforced-space notation, Conjecture 4.22, and the
extension problem found the source thesis and related reinforced-space work,
but no paper stating this equality.  Lu and Yao prove the reflection
extension theorem used here, but do not discuss reinforced face weights or
Scharf's conjecture.  Novelty is therefore plausible, not certified.

\begin{thebibliography}{9}
\bibitem{Scharf}
B.~Scharf,
\emph{Wavelets in function spaces on cellular domains},
arXiv:1302.3751 (2013).

\bibitem{LuYao}
H.~Lu and L.~Yao,
\emph{On Seeley-type universal extension operators for the upper half
space}, arXiv:2211.15567v2 (2024), especially Theorem~23 and
equations (29)--(30).
\end{thebibliography}

\end{document}
